尤其需要指出的是，金融市场极其复杂，它可能会受到自然灾害、非平稳数据分布、噪声信号、以及宏观经济或意外事件导致的突然政权更迭、某些类型的虚假新闻等因素的影响，这些相互关联的因素每天都在发生。这使得预测价格变得非常困难，但对于股票投资者和交易员来说，预测价格仍然至关重要。本文通过结合稳健的算法策略、自适应技术指标优化、深度学习和深度强化学习架构来解决这些问题，这些架构可以在不断变化的市场条件下进行泛化，从而预测市场上的股票价格，特别是针对英伟达公司 (NVDA) 和标准普尔 500 指数 (GSPC)，并使用截至 2025 年 3 月 5 日的可用指数数据。

简而言之，对一组 20 种基于 PineScript 的交易策略进行了统计回测，这些策略称为“FinDRL”，源自趋势跟踪 (TF)、突破检测 (BO)、均值回归 (MR) 和带有随机 RSI 屏幕的成交量分析。 FinDRL - Momentum、FinDRL - Williams \%R 和 FinDRL - Volume Spike 是表现最出色的指标之一，拥有可观的利润因子（例如 2.469），并展示了算法杠杆在复杂市场条件下的价值。其他指标则存在局限性，强调了选择性应用的必要性。

基于这些初步结果，应用了基于 Optuna Bayesian 优化的优化框架来提高这些指标的有效性。通过系统参数调整，IndicatorOptimizer 类将曾经被认为无效的指标（如 Parabolic SAR）转变为有利可图的交易策略，在利润和回撤减少方面实现了多个级别的性能提升。验证是在高保真交易环境中进行的，使该练习与现实世界的投资组合管理相关。

然后，基于 GRU、CNN 和 LSTM 架构训练和优化更复杂的 DL 模型，以进一步提高预测能力。这些模型被纳入一个动态回测框架，该框架部署了集合预测和优化的回溯期（10-120 天），这些模型始终远远超过传统基线。我们使用了包括最大回撤和夏普比率在内的指标，发现预测准确性和交易能力都有显著提高。

结果促成了一个名为 FinDRL 的复杂 DRL 系统的建立，该系统封装了贝叶斯优化的深度神经网络 (DNN)（用作特征提取器）和深度 Q 网络 (DQN)，以便在现实交易环境 (TradingEnv) 中做出可操作的决策。该环境包括 OHLCV 数据、简单移动平均线 (SMA) 指标和交易成本。实证验证显示投资组合增长稳健（15-20\%），风险调整后表现良好（夏普比率：1.2）。可视化 Quest 可视化提供全面的可视化（例如股权曲线、交易信号、绘图等），以便清晰地回顾分析和比较策略的表现。

总体而言，这项研究提出了一种模块化、可组合和可扩展的架构，将传统技术指标与先进的 DL/DRL 一致的表示形式相结合，从而为市场参与者和研究人员提供基础工具和见解。